\slajdid{01-s030}
\begin{frame}[label=01-s030]{Aktivan logički nivo}
\gusttekst
Naviknuti da razmišljamo u „pozitivnoj“ logici

\begin{description}
\item[I logičko kolo] – daće logičku jedinicu na izlazu ako su na svim ulazima logičke jedinice\\
\alert{i} na prvom \alert{i} na drugom \alert{i} …
\item[ILI logičko kolo] – daće logičku jedinicu na izlazu ako je bar na jednom ulazu logička jedinice\\
\alert{ili} na prvom \alert{ili} na drugom \alert{ili} …
\end{description}

Znači možemo da kažemo da logička kola I i ILI izvršavaju svoju logičku funkciju i kada su zadovoljeni logički uslovi na izlaz će biti logička jedinica

\alert{Znači možemo da kažemo da logička kola I i ILI izvršavaju svoju logičku funkciju i kada su zadovoljeni logički uslovi na izlaz će biti aktivan nivo logičke jedinica}
\end{frame}
